\section{Contexto}

El Trabajo Práctico N.º~3 de la materia Arquitectura de Computadoras
propone crear un procesador \tech{RISC-V} de 32 bits con un pipeline
de cinco etapas, una unidad de debug y un enlace \tech{UART} que
permita operarlo desde una computadora. El resultado debe ser un
sistema que sea capaz de recibir un programa escrito en
\tech{assembler}, lo traduzca a código máquina, lo cargue en la
memoria de instrucciones, lo pueda ejecutar en modo continuo o paso a
paso y, finalmente, devuelva una observación coherente del estado interno.

\section{Alcance del trabajo}

El sistema comprende tres partes principales:

\begin{enumerate}
  \item un procesador \tech{RISC-V} de 32 bits con pipeline de cinco etapas;
  \item una unidad de carga, control de ejecución y debug conectada
    con una PC por \tech{UART};
  \item una aplicación de PC para editar, traducir, cargar, ejecutar y observar
    programas.
\end{enumerate}

El procesador implementa un subconjunto de \tech{RV32I} de 32
instrucciones y la instrucción \code{halt}. La versión del ISA
adoptada es User-Level ISA V2.2. La FPGA de referencia es una
\tech{Basys 3} (\tech{XC7A35T-1CPG236C}) con un dominio funcional de
\SI{50}{\mega\hertz}. El perfil de memoria de referencia utiliza
\SI{1024}{\byte} de IMEM y \SI{2048}{\byte} de DMEM.

\section{Metodología de desarrollo}

El desarrollo siguió una secuencia de especificación, implementación
y verificación. En la primera etapa se identificaron los requisitos,
se registraron las decisiones de arquitectura y se definieron las
interfaces entre bloques. Esa documentación permitió implementar cada
módulo contra un contrato explícito y construir bancos de prueba
antes de conectar los bloques entre sí.

La implementación se organizó por entregas funcionales. Cada entrega
reunió un conjunto coherente de módulos, sus pruebas y la evidencia
de ejecución correspondiente. Las etapas del procesador se
verificaron primero de forma aislada y después en los niveles de
integración local y CPU. La comunicación, el control de sesión y la
plataforma se incorporaron sobre esa base, conservando las interfaces
y los identificadores de las pruebas.

La verificación combina simulaciones, revisiones estructurales y
análisis de síntesis e implementación. Las pruebas funcionales se
ejecutan con Icarus Verilog; el análisis del RTL se realiza con
Verilator y Verible; la implementación física se evaluó con Vivado.
Los resultados se registran por caso, nivel y configuración de memoria.

La especificación completa de la arquitectura se encuentra en
\doc{architecture/system_overview.md}{el diseño FPGA consolidado}.
Los requisitos y sus criterios de aceptación se detallan en
\doc{requirements.md}{la especificación de requisitos}. La
organización de los bloques y sus interfaces se documenta en
\doc{architecture/module_inventory.md}{el inventario de módulos} y en
\doc{interfaces.md}{el contrato de interfaces}.
